\section{Introduction}\label{sec:intro}

\subsection{Traffic Forecasting in Urban Mobility}\label{sec:intro-background}

Road transport is a major mode of movement in many cities, and the demand
placed on urban networks continues to grow with population and economic
activity. Congestion imposes costs on travelers, freight operators and the
environment, while the physical capacity of the road network is difficult to
expand in built-up areas. Operators therefore rely on traffic management to use
existing capacity efficiently. Short-term traffic forecasting is a core part of
this task: it estimates flow, occupancy or speed over the next 15 to 60 minutes,
giving signal controllers, route guidance systems, incident response teams and
network planners time to act before congestion forms.

Forecasting methods have developed alongside the sensing infrastructure that
supplies them. Loop detectors, radar and probe vehicles provide regular time
series at fixed locations, and early methods modeled each location
independently. As sensor networks expanded, spatio-temporal graph neural
networks became a prominent approach: they represent sensors as nodes in a graph
and learn both the spatial correlation between neighboring locations and the
temporal dynamics of the series \parencite{yu2018stgcn,wu2019gwnet}. Benchmarks such as LibCity
provide common data splits and implementations, making systematic comparison
possible \parencite{wang2021libcity,jiang2023libcity}.

Camera networks add a complementary view. Fixed roadside cameras cover wide
areas and record individual vehicles, and one-stage detectors can recover
vehicle positions and classes from each frame \parencite{ultralytics2026yolo}.
The resulting counts describe what is visible at the camera, while calibrated
flow, occupancy and speed series describe the aggregate traffic state used by
forecasting models. Combining the two views is a practical route toward
monitoring systems that both observe and anticipate traffic.

\subsection{Problem Statement and Project Scope}\label{sec:intro-objectives}

Singapore operates a dense network of fixed roadside cameras that provide
continuous visual coverage of its road network. The LTA collection used in this
project contains 15,316 images from eight cameras over two weeks, sampled at
ten-minute intervals. The images support vehicle recognition and manual
monitoring. They do not include calibrated flow, occupancy or speed at the
camera locations, and the ten-minute sampling interval gives a discontinuous
view of passing traffic. A forecasting model, by contrast, requires a regular
series of aggregate measurements. Public benchmarks provide such series, but for
different road networks.

The project addresses this mismatch by building a prototype with three
components: a vehicle recognition module, a traffic forecasting module and the
Traffic Monitoring and Forecasting Workbench. The detection module is trained on
MIO-TCD and jointly fine-tuned on LTA images with teacher pseudo-labels. The forecasting module is trained and
evaluated on PEMSD8 under a common protocol that compares eight models. The
workbench presents the outputs of both modules and labels their data sources. The
design keeps the two data domains explicit and defines the interface through
which a calibrated local series would connect them.

The project pursues three questions:

\begin{itemize}
  \item How well does a detector trained on MIO-TCD transfer to LTA images,
        measured against teacher labels and against a human-labeled LTA test
        set?
  \item How do representative forecasting models compare on PEMSD8 under a
        common protocol, and how does the custom model perform against them?
  \item How can a workbench present both outputs while keeping their data sources
        explicit?
\end{itemize}

The scope is limited to four vehicle classes in still images. Forecasting
focuses on traffic flow, with occupancy and speed as secondary analyses. The
project is a time-bounded prototype and pipeline design. It establishes the
processing chain, evaluation protocol and interface. An operational Singapore
forecasting system would additionally require calibrated sensing and longer
data collection. Nighttime LTA frames are a particular constraint: low light and
small vehicle scale make annotation difficult even for humans.

\subsection{Contributions}\label{sec:intro-contributions}

\begin{itemize}
  \item An LTA image set of 15,316 frames from eight Singapore cameras over two
        weeks, with a documented sampling and teacher pseudo-labeling pipeline,
        and a human-labeled test set of 63 week-2 images (1,197 boxes) frozen
        before evaluation.
  \item A leakage-aware MIO-TCD evaluation protocol based on scene groups, and
        a joint fine-tuning procedure that trains a COCO-pretrained YOLO26m on
        MIO-TCD and LTA images with teacher pseudo-labels; the model is applied
        to the full LTA collection.
  \item A per-camera false-positive filter calibrated by hand for all eight
        cameras, and a congestion-level estimate for all 15,316 LTA frames that
        combines the detector, a classical vision method and a zero-shot
        vision-language model; both are evaluated against human labels, with
        method selection and testing kept on separate label sets.
  \item A common-protocol benchmark of eight forecasting models on PEMSD8,
        covering native and unified nine-channel inputs, with five-seed paired
        checks for the top four models.
  \item MultiHeadSTGCN, an STGCN-based model with direct multi-step output,
        calendar inputs and a selectable spatial operator, examined in a staged
        ablation and against the benchmark models.
  \item A web workbench that integrates overview, camera monitoring and
        forecast views, with historical replay, camera-reliability indicators
        and a threshold-based risk ranking, together with an integration design
        for a calibrated local series.
\end{itemize}
